\documentclass[12pt]{article}
\usepackage{amsmath}
\usepackage{amssymb}

\begin{document}

\noindent\textbf{Question:} How does the computational cost of the particle filter scale with the number of particles and the number of dimensions in the state vector? Why can large dimensionality be a problem for particle filters in practice?

\vspace{0.5cm}

\noindent\textbf{Scaling with Number of Particles ($N$):}

Computational cost scales \textbf{linearly} with $N$. Each iteration requires $O(N)$ operations for prediction, $O(N \cdot M)$ for measurement updates (where $M$ is the number of measurements), and $O(N)$ for resampling.

\vspace{0.5cm}

\noindent\textbf{Scaling with State Dimensionality ($d$):}

The number of particles required for adequate state space coverage scales \textbf{exponentially} with dimensionality:
\[
N_{\text{required}} \propto \epsilon^{-d}
\]
where $\epsilon$ is the resolution per dimension.

\vspace{0.5cm}

\noindent\textbf{Why High Dimensionality is Problematic:}

\begin{enumerate}
    \item \textbf{Exponential particle growth}: A 3D state $(x, y, \theta)$ might need $10^3$ particles, while a 6D state needs $10^6$ particles for comparable coverage.
    
    \item \textbf{Weight degeneracy}: In high dimensions, the probability that any particle is close to the true state in \textit{all} dimensions decreases exponentially. Most particles receive negligible weights, causing effective sample size collapse:
    \[
    N_{\text{eff}} = \frac{1}{\sum_{i=1}^{N} w_i^2} \ll N
    \]
    
    \item \textbf{Computational intractability}: Memory and computation scale as $O(\epsilon^{-d} \cdot d)$, making high-dimensional problems ($d > 4$) impractical with standard particle filters.
\end{enumerate}

\end{document}

\end{document}